\documentclass[parskip=full]{scrartcl}
\usepackage{geometry}
\geometry{a4paper, left=2.5cm, right=2cm, top=3cm, bottom=3cm}
\usepackage[utf8]{inputenc} % use utf8 file encoding for TeX sources
\usepackage[T1]{fontenc}    % avoid garbled Unicode text in pdf
\usepackage[german]{babel}  % german hyphenation, quotes, etc
\usepackage{hyperref}       % detailed hyperlink/pdf configuration
\hypersetup{                % ‘texdoc hyperref‘ for options
pdftitle={Echtzeit-Computergrafik in der Spieleentwicklung},%
bookmarks=true,
}

\usepackage{rotating}
\usepackage{newclude}
\usepackage{typearea}
\newcommand\invisiblesection[1]{%
  \refstepcounter{section}%
  \addcontentsline{toc}{section}{\protect\numberline{\thesection}#1}%
  \sectionmark{#1}}


\usepackage{listings}

\usepackage{xcolor}
\usepackage{graphicx}       % provides commands for including figures
\usepackage{csquotes}       % provides \enquote{} macro for "quotes"
\usepackage[nonumberlist]{glossaries}     % provides glossary commands
\usepackage{enumitem}
\usepackage{epsf,epsfig,eepic}
\usepackage[headsepline, footsepline]{scrpage2}


%Design für Methoden-/Attributbeschreibungen
\usepackage{framed}
\usepackage{color}
\definecolor{gray}{RGB}{200,200,200}
\renewenvironment{leftbar}[1][\hsize]
{%
\def\FrameCommand
{%
{\color{gray}\vrule width 1pt}%     
}%
\MakeFramed{\hsize#1\advance\hsize-\width\FrameRestore}%
}
{\endMakeFramed}

\definecolor{eclipseStrings}{HTML}{B71C1C}
\definecolor{eclipseKeywords}{HTML}{0D47A1}
\definecolor{background}{HTML}{ECEFF1}
\colorlet{numb}{magenta!60!black}

\lstdefinelanguage{json}{
    basicstyle=\normalfont\ttfamily,
    commentstyle=\color{eclipseStrings}, % style of comment
    stringstyle=\color{eclipseKeywords}, % style of strings
    numbers=left,
    numberstyle=\scriptsize,
    stepnumber=1,
    numbersep=8pt,
    showstringspaces=false,
    breaklines=true,
    frame=lines,
    backgroundcolor=\color{background}, %only if you like
    string=[s]{"}{"},
    comment=[l]{:\ "},
    morecomment=[l]{:"},
    literate=
        *{0}{{{\color{numb}0}}}{1}
         {1}{{{\color{numb}1}}}{1}
         {2}{{{\color{numb}2}}}{1}
         {3}{{{\color{numb}3}}}{1}
         {4}{{{\color{numb}4}}}{1}
         {5}{{{\color{numb}5}}}{1}
         {6}{{{\color{numb}6}}}{1}
         {7}{{{\color{numb}7}}}{1}
         {8}{{{\color{numb}8}}}{1}
         {9}{{{\color{numb}9}}}{1}
}

%kopf und fusszeilen
\pagestyle{scrheadings}
\clearscrheadfoot

%kopfzeile
\ihead{\codename}
\ohead{\headmark}
\automark{section}

%fusszeile
\ofoot{\pagemark}
\ifoot{\codename}

\makenoidxglossaries
%
% % Glossareinträge
%

%Kommandos
\newcommand{\codename}{Valaris}


\title{Implementierungsdokument}
\author{Sidney Hansen, Jonas Heinle, Frederik Lingg, Lukas Schölch, Artur Wesner}

\begin{document}

	\begin{titlepage}
		\centering
		{\scshape\LARGE Implementierungsdokument\par}
		\vspace{1cm}
		{\scshape\Large Echtzeit Computergrafik in der Spieleentwicklung \par}
		\vspace{1cm}
		{\huge\bfseries \codename \par}
		\vspace{1cm}
		\includegraphics[width=.5\linewidth]{./Bilder/valaris_t_4096.png}
		\par
		{\vspace{1cm}}
		{\Large\itshape Autoren \\}
		{\Large\itshape Sidney Hansen, Jonas Heinle,\\}
		{\Large\itshape Frederik Lingg, Lukas Schölch, Artur Wesner\par}
		
		\vfill
		Projektbetreuer\par
		Alisa Jung \& Tobias Rapp
		
		\vfill
		
		% Bottom of the page
		{\large \today\par}
	\end{titlepage}

	\tableofcontents 
	\pagebreak
	
	
	\pagebreak
	
	\section{Einleitung}

	Für die Implementierung von Valaris wurde die Gliederung, welche bereits im Pflichtenheft und Entwurf eingehalten wurde,
	weiterhin verwendet. Somit wurden die Module InGame-Grafik, InGame-Interface, Assets, Generierung und GUI zunächst unabhängig voneinander implementiert
	und später zusammen gefügt.
	
	Jedoch wurde aufgrund von ungleichem Arbeitsumfang der Module die Arbeitsaufteilung der Module weniger streng gehalten.
	Zusätzlich ging Zeit verloren, da das zusammenführen der einzelnen Module zum Teil länger dauerte als ursprünglich vorgesehen.
	Ausserdem wurde die Implementierung der Schnittstelle zum anderen Team mehrfach überarbeitet, um komplexität im Simulationsthread zu verringern.
	Die Implementierung der Shader für Physik basierte Materialien erwies sich ebenfalls als komplizierter als gedacht, da es häufig zu Rechenfehlern
	aufgrund der Ungenauigkeit von Gleitkomma-Zahlen kam.

	Das Resultat der Implementierung ist nun eine Anwendung, welche die gesammte Menüführung besitzt, Strecken generieren und vorführen kann.
	Die eigentliche Simulation des Spiels wurde bisher nicht eingegliedert.

	Im Folgenden wird im Detail erleutert, welche Kriterien umgesetzt, welche Aspekte des Entwurfs geändert und welche Unit-Tests verwendet wurden.
	Das Dokument schließt mit dem Vergleich des geplanten und dem tatsächlichen Ablauf der Implementierung anhand von Gantt-Diagrammen ab.

	\pagebreak

	\section{Features}
	\subsection{InGame - Grafik}
\subsubsection{Musskriterien}

\paragraph{Implementiert}
\begin{itemize}
    \item Darstellung physikalisch akkurater Materialien
    \item Darstellung metallener bzw. dielektrischer Materialien
		\item Rendern von undurchsichtigen Objekten
		\item Viele dynamische Lichtquellen gleichzeitig verarbeiten
		\item Schatten von dynamischen Lichtquellen
		\item Effekte
		\begin{itemize}
            \item verschiedene Partikeleffekte (z.B. Explosionen)
        \end{itemize}
        \item Frei fliegende Kamera zum Ansehen der Strecke
\end{itemize}

\subsubsection{Wunschkriterien}

\paragraph{Implementiert}
\begin{itemize}
    \item Animationen an Fahrzeug
    \begin{itemize}
        \item Animieren der Düsen des Fahrzeuges
    \end{itemize}
    \item Leuchtende Oberflächen
\end{itemize}

\paragraph{Nicht Implementiert}
\begin{itemize}
    \item Postprocessing (z.B. MotionBlur)
	\item Animationen an Spielfigur und Fahrzeug
    \begin{itemize}
	\item Klauen der Energie am vorderen Fahrzeug
	\end{itemize}
	\item Verschiedene Fahrzeuge
    \item HDR
	\item FXAA
	\item Soundeffekte
\end{itemize}
	\pagebreak
	\subsection{InGame - Interface}
Im Pflichtenheft wurden zu diesem Modul keine speziellen Anforderungen genannt, daher
werden im Folgenden alle Implementierten-Features als Musskriterien gewertet.

\subsubsection{Musskriterien}
\paragraph{Implementiert}
\begin{itemize}
    \item Übergabe der Daten in voneinander unabhängigen Paketen
    \item Minimale warte Zeiten bei der Übergabe der Pakete
    \item Dynamische Anzahl von Spiel-Objekten in der Szene
    \item Dynamische Anzahl von Eigenschaften eines einzelnen Spiel-Objekts
    \item Übergabe von Spiel-Events (zum Beispiel ende des Rennens) durch das Interface

    \item Eigenschaften für Position, Ranking, Energielevel, Camera, Türstatus, Geschwindigkeit

    \item Laden der Modelle für Spiel-Objekte
    \item Interpretieren der Eigenschaften der SpielObjekte zur korrekten Darstellung in der Szene
\end{itemize}

\paragraph{Nicht Implementiert}
Aufgrund mangelnder Definition von Anforderungen an das Interface gibt es hier keine nicht implementierten
Musskriterien
	\pagebreak
	\subsection{Assets}
\subsubsection{Musskriterien}

\paragraph{Implementiert}
\begin{itemize}
    \item VerbindungsTunnel der Raumstation
    \item Dekorationen für Raumstationselemente
    \item Kuppeln mit planetarischer Umgebing
    \item Hover-Cart
\end{itemize}

\paragraph{Nicht Implementiert}
\begin{itemize}
    \item Energie-Kugeln auf Hover-Cart
    \item Fahrer für Hover-Cart
\end{itemize}

\subsubsection{Wunschkriterien}

\paragraph{Nicht Implementiert}
\begin{itemize}
    \item Verschiedene Verbindungstunnel-Typen
    \item Verschiedene Fahrzeug-Modelle
    \item Animationen für Fahrer
\end{itemize}
	\pagebreak
	\subsection{Generierung}
\subsubsection{Musskriterien}

\paragraph{Implementiert}
\begin{itemize}
    \item Zusammenhängende Strecke mit Start und Ziel
    \item Checkpoints in sinnvollen Abständen
    \item Generierung von unterschiedlichen Streckenführungen
    \item Deterministisch vom Seed abhängig
    \item Keine unbefahrbaren Streckenabschnitte
    \item Definieren von Extremparametern (Bsp. Steigung, Schärfe der Kurven, ...)
    \item Konform zum Streckenmodell des anderen Teams
    \item (Generierung von Streckenumgebung)
\end{itemize}

\paragraph{Nicht Implementiert}
\begin{itemize}
    \item Assets als Streckenumgebung in Kuppeln
\end{itemize}

\subsubsection{Wunschkriterien}

\paragraph{Implementiert}
\begin{itemize}
    \item Verschiedene Biome
		
\end{itemize}

\paragraph{Nicht Implementiert}
\begin{itemize}
    \item Verschiedene befahrbare Untergründe
	\item Einfügen von Hindernissen auf die Rennbahn
    \item Dynamische Strecken- bzw. Umgebungselemente
\end{itemize}
	\pagebreak
	\subsection{GUI}
\subsubsection{Musskriterien}

\paragraph{Implementiert}
\begin{itemize}
    \item Allgemeines Start-Menü
	\item Einstellungsfenster
    \begin{itemize}
        \item Diverse Möglichkeiten zur Konfiguration des Spielverhaltens
        \item Diverse Möglichkeiten zur Konfiguration der Grafik
    \end{itemize}
    \item Vor dem Spiel:
    \begin{itemize}
        \item Konfiguration des Seed
        \item Auswählen eines gespeicherten Seed
        \item Definition der Anzahl an optionaler KI
    \end{itemize}
    \item Ein während des Spiels aufrufbares Pause-Menü mit den Funktionen:
    \begin{itemize}
        \item Zurückkehren zum Spiel
        \item Zurückkehren zum Start-Menü
        \item Öffnen der Einstellungen
        \item Speichern des aktuellen Seed
    \end{itemize}
    \item Nach dem Spiel:
    \begin{itemize}
        \item Möglichkeit zur Speicherung des aktuellen Seed
        \item Eine Ergebnisanzeige mit Visualisierung der Endposition und der gefahrenen Zeit eines jeden Spielers
        \item Möglichkeit die Strecke erneut zu fahren
        \item Zurückkehren zum Start-Menü
    \end{itemize}
    \pagebreak
    \item Anzeige eines Head-up-Displays(HUD) während dem Spiel
    \begin{itemize}
        \item Aktuell gefahrene Zeit
        \item Visualisierung des aufgesammelten Powerups
        \item Anzeige der aktuellen Ranglistenplatzierung
    \end{itemize}
\end{itemize}

\paragraph{Nicht Implementiert}
\begin{itemize}
    \item Vor dem Spiel:
    \begin{itemize}
        \item Auswählen des Spielmodus
    \end{itemize}
    \item Anzeige eines Head-up-Displays(HUD) während dem Spiel
    \begin{itemize}
        \item Mini-Map in 2D inklusive Visualisierung der Fahrer
        \item Anzeige der Rundenzahl
        \item Anzeige der Geschwindigkeit des Fahrers
    \end{itemize}
\end{itemize}

\subsubsection{Wunschkriterien}

\paragraph{Implementiert}

\paragraph{Nicht Implementiert}
\begin{itemize}
    \item Konfigurationseinstellung für den Ton
    \item Fahrzeugauswahl
    \item Beispielvorschau eines Seed
    \item Abstimmmöglichkeit unter Spielern, ob eine Strecke erneut gespielt werden soll
\end{itemize}
	\pagebreak

	\section{Änderungen}
	\subsection{InGame - Grafik}

\subsubsection{Neue Features}
\begin{itemize}
    \item \textit{Testfeature}
        \begin{leftbar}[0.9\linewidth]
            Beschreibung des Features
        \end{leftbar}
\end{itemize}

\subsubsection{Änderungen existierender Komponenten}
\begin{itemize}
    \item \textit{Testfeature}
        \begin{leftbar}[0.9\linewidth]
            Beschreibung der Änderung und evtl Begründung
        \end{leftbar}
\end{itemize}

\subsubsection{Nicht implementiert}
\begin{itemize}
    \item \textit{TestNotExistingfeature}
        \begin{leftbar}[0.9\linewidth]
            Alle Methoden die beim Partikeleffekt abhängig von tpf sind wurden nicht
            mit einen JUnit test getestet. Da Ergebnisse abhängig von Framerate macht
            ein klassischer JUnit Test keinen Sinn. Diese wurden allerdings ausführlich
            grafisch auf Funktionalität getestet. 
        \end{leftbar}
\end{itemize}
	\pagebreak
	\subsection{InGame - Interface}

\subsubsection{Neue Features}
\subsubsection{Änderungen existierender Komponenten}
\begin{itemize}
    \item \textit{Tick ist Threadsafe}
        \begin{leftbar}[0.9\linewidth]
            Die Klasse Tick wurde durch das hinzufügen von Synchronisationsmechanismen robuster gemacht.
        \end{leftbar}
\end{itemize}
\subsubsection{Nicht implementiert}
	\pagebreak
	\subsection{Assets}

\subsubsection{Neue Features}
\begin{itemize}
    \item \textit{Metadaten in AssetDescriptoren}
        \begin{leftbar}[0.9\linewidth]
            Es ist nun möglich Key-Value Paare im AssetDescriptor zu definieren und diese
            zur Laufzeit durch das AssetControl aus zu lesen
        \end{leftbar}
\end{itemize}

\subsubsection{Änderungen existierender Komponenten}
\begin{itemize}
    \item \textit{Laden von Controls}
        \begin{leftbar}[0.9\linewidth]
            Das AssetControl kann nun ausschließlich Controls vom Typ AbstractAssetControl laden,
            ist dadurch aber in der Lage geladene Controls über das zerstören ihres Parents zu benachrichtigen.
        \end{leftbar}
\end{itemize}
\subsubsection{Nicht implementiert}
	\pagebreak
	\subsection{Generierung}

\subsubsection{Neue Features}
\begin{itemize}
    \item \textit{Bepflanzung der Kuppeln}
        \begin{leftbar}[0.9\linewidth]
            Es können nun für ein bestimmtes Biom eine Menge von Bäumen(bzw. ähnliche Pflanzen) definiert werden,
            welche dann in einer entsprechenden Kuppel platziert werden. Bei der Platzierung werden die Objekte
            anhand einer, im \textit{DomeAssetGenerator} definierten, Dichtefunktion verschoben, wie stark 
            kann durch einen Parameter pro Biom bestimmt werden.
            Desweiteren kann man pro Biom unter anderem die Asset-Skalierung und Farben der Blätter festlegen.
        \end{leftbar}

    \item \textit{Assets als Kuppeldach}
        \begin{leftbar}[0.9\linewidth]
            Es ist nun möglich Assets als Dach für eine Kuppel zu verwenden. Es können auch verschiedene Assets
            definiert werden, aus denen dann ein zufälliges(abhängig vom Seed) ausgesucht wird.
        \end{leftbar}
\end{itemize}

\subsubsection{Änderungen existierender Komponenten}
\begin{itemize}
    \item \textit{Methode \textit{generateScenegraph()} parallelisiert}
        \begin{leftbar}[0.9\linewidth]
            \textit{GenerateScenegraph()} der Klasse Map, ruft nun die \textit{generateScenegraph()} Methoden
            der ISceneItem parallelisiert auf und beschleunigt so den Generierungsprozess deutlich.
        \end{leftbar}
    
    \item \textit{'Colring' von Landschaften verschönert}
        \begin{leftbar}[0.9\linewidth]
            Es entstehen nun keine ungewollten (bunte) Dreiecksstrukturen mehr, da das festlegen der 'Oberflächen'
            an einem Punkt optimiert wurde.
        \end{leftbar}

    \item \textit{Verbesserte Räume}
        \begin{leftbar}[0.9\linewidth]
            Ein weiterer Raumtyp wurde hinzugefügt, sowie die Asset-Platzierung in räumen optimiert.
        \end{leftbar}

    \item \textit{Unterschiedliche Strecken}
        \begin{leftbar}[0.9\linewidth]
            Unterschiedliches Straßen Assets in Tunneln, Kuppeln und Räumen.
        \end{leftbar}

    \item \textit{Neue Biome}
        \begin{leftbar}[0.9\linewidth]
            Es wurden 4 fertige Biome definiert, die sich in den Punkten Untergrund, Farben und platzierten Objekten
            unterscheiden.
        \end{leftbar}

    \item \textit{Kuppeldach wurde durch Asset ersetzt}
        \begin{leftbar}[0.9\linewidth]
            Das Dach einer Kuppel besteht nun aus einem, in Blender erstelltem, Asset.
        \end{leftbar}

    \item \textit{Definition von 3 Kartentypen}
        \begin{leftbar}[0.9\linewidth]
            Es wurden ein kleines, ein mittleres und ein großes Biom definiert, die sich in der Anzahl der Kuppeln, 
            Tunneln und Räumen unterscheiden. Ebenso sind di möglichen Größen eines solchen \textit{ISceneItems} unterschiedlich.
        \end{leftbar}
\end{itemize} 
	\pagebreak
	\subsection{GUI}

\subsubsection{IConfig}
\begin{itemize}
    \item \textit{Umbenannt zu ConfigManager und kein Interface mehr}
        \begin{leftbar}[0.9\linewidth]
            Hier lag ein Entwurfsfehler vor. Diese Klasse stellt nun Funktionen zum Lesen und 
            Schreiben einer Datei bereit. Sie enthält für jede Methode einen Pfad, wo sich 
            die Datei befinden sollte.
        \end{leftbar}
\end{itemize}

\subsubsection{InternalGameConfig, SeedConfig und GraphicsConfig}
\begin{itemize}
    \item \textit{Implementiert nicht mehr IConfig}
        \begin{leftbar}[0.9\linewidth]
            Da die Klasse statische Methoden enthält, implementiert sie nicht mehr das 
            Interface IConfig.
        \end{leftbar}
    \item \textit{Methode getConfig() hinzugefügt}
        \begin{leftbar}[0.9\linewidth]
            Gibt die Konfiguration zurück und ruft bei die getConfig()
            Methode vom ConfigManager auf.
        \end{leftbar}
    \item \textit{Methode saveConfig() hinzugefügt}
        \begin{leftbar}[0.9\linewidth]
            Speichert die Konfiguration und ruft bei die 
            saveConfig() Methode vom ConfigManager auf.
        \end{leftbar}
    \item \textit{Methode reset() hinzugefügt}
        \begin{leftbar}[0.9\linewidth]
            Setzt die Config zurück durch Überschreiben mit der Default-Datei.
        \end{leftbar}
\end{itemize}

\subsubsection{OwnScreenController}
\begin{itemize}
    \item \textit{Klasse hinzugefügt}
        \begin{leftbar}[0.9\linewidth]
            OwnScreenController implementiert von ScreenController und wird von jedem Menü-
            Bildschirm geerbt. Es enthält Funktionen und Attribute, die jeder Bildschirm 
            braucht und reduziert somit Code-Duplizierung. Jedem OwnScreenController wird 
            außerdem eine Referenz auf den NiftyAppState und die anderen OwnScreenController
            mitgegeben, so dass diese leicht aufrufbar sind.
        \end{leftbar}
\end{itemize}

\subsubsection{NiftyAppState}
\begin{itemize}
    \item \textit{Gesamtes Entwurfskonzept der Klasse geändert}
        \begin{leftbar}[0.9\linewidth]
            Ein NiftyAppState ist nun eine AppState, welcher alle die ihm mitgegebenen 
            OwnScreenController in einer Nifty Instanz initialisiert und speichert.
            Die OwnScreenController innerhalb eines NiftyAppStates haben Zugriff auf 
            einander und können so leicht mit nifty.goToScreen() aufgerufen werden.
            Ist dies nicht möglich, so erstellt ein OwnScreenController ein neues 
            NiftyAppState-Objekt mit den jeweiligen OwnScreenControllern und fügt dieses
            an den AppStateManager an.
            Wird der NiftyAppState vom AppStateManager entfernt, so entfernt dieser alle 
            initialisierten Sachen wie z.B den Processor für das NiftyDisplay.
        \end{leftbar}
\end{itemize}

\subsubsection{SplashScreen}
\begin{itemize}
    \item \textit{Klasse wurde hinzugefügt}
        \begin{leftbar}[0.9\linewidth]
            Zeigt einen SplashScreen bei Start des Programms an.
        \end{leftbar}
\end{itemize}

\subsubsection{MainmenuScreen}
\begin{itemize}
    \item \textit{Implementiert nicht mehr ScreenController, erbt dafür von OwnScreenController}
        \begin{leftbar}[0.9\linewidth]
            Es wurde eine eigene ScreenController Klasse hinzugefügt.
        \end{leftbar}
    \item \textit{Methode openMultiplayerMode() entfernt}
        \begin{leftbar}[0.9\linewidth]
            Wurde nicht mehr benötigt, da die Mehrspieler-Funktion inklusive des 
            Lobby Bildschirmes entfernt wurde.
        \end{leftbar}
\end{itemize}

\subsubsection{AbstractEditGameSettings}
\begin{itemize}
    \item \textit{Klasse entfernt}
        \begin{leftbar}[0.9\linewidth]
            Diese Klasse existierte ursprünglich, da der Lobby Bildschirm und der 
            Einzelspielerbildschirm viele gleiche Methoden brauchten. Da der Lobby 
            Bildschirm entfernt wurde, wurden alle Funktionen hierraus in die Klasse 
            SinglePlayerModeScreen verschoben.
        \end{leftbar}
\end{itemize}

\subsubsection{LobbyScreen}
\begin{itemize}
    \item \textit{Klasse entfernt}
        \begin{leftbar}[0.9\linewidth]
            Wurde nicht mehr benötigt, da die Mehrspieler-Funktion inklusive des 
            Lobby Bildschirmes entfernt wurde.
        \end{leftbar}
\end{itemize}

\subsubsection{SinglePlayerModeScreen}
\begin{itemize}
    \item \textit{Implementiert nicht mehr ScreenController, erbt dafür von OwnScreenController}
        \begin{leftbar}[0.9\linewidth]
            Es wurde eine eigene ScreenController Klasse hinzugefügt.
        \end{leftbar}
    \pagebreak
    \item \textit{Methoden aus AbstractEditGameSettings hier eingefügt}
        \begin{leftbar}[0.9\linewidth]
            AbstractEditGameSettings wurde entfernt und dementsprechend sind dessen Methoden 
            hier eingesetzt worden.
        \end{leftbar}
    \item \textit{Methode editCurrentSeed(String, DropDownSelectionChangedEvent) hinzugefügt}
        \begin{leftbar}[0.9\linewidth]
            Wird aufgerufen sobald eine Auswahl im DropDown-Menü getroffen wird.
        \end{leftbar}
    \item \textit{Methode editCurrentSeedNumber(String, TextFieldChangedEvent) hinzugefügt}
        \begin{leftbar}[0.9\linewidth]
            Wird aufgerufen sobald die Seed Nummer bearbeitet wird.
        \end{leftbar}
    \item \textit{Methode editKICount(String, TextFieldChangedEvent) hinzugefügt}
        \begin{leftbar}[0.9\linewidth]
            Wird aufgerufen sobald die KI Anzahl bearbeitet wird.
        \end{leftbar}
    \item \textit{Methode saveSeedName() hinzugefügt}
        \begin{leftbar}[0.9\linewidth]
            Wird aufgerufen sobald der Seed Name mit Enter bestätigt, oder der Speicher-Button
            rechts daneben geklickt wird.
        \end{leftbar}
    \item \textit{Methode editSeednameTextfieldVisibility() hinzugefügt}
        \begin{leftbar}[0.9\linewidth]
            Wird aufgerufen sobald der Bearbeiten-Button neben dem DropDown-Menü gedrückt wird.
        \end{leftbar}
    \item \textit{Methode saveSeedNumber() hinzugefügt}
        \begin{leftbar}[0.9\linewidth]
            Wird aufgerufen sobald die Seed Nummer mit Enter bestätigt, oder der Speicher-Button
            rechts daneben geklickt wird.
        \end{leftbar}
\end{itemize}

\subsubsection{SettingsmenuScreen}
\begin{itemize}
    \item \textit{Implementiert nicht mehr ScreenController, erbt dafür von OwnScreenController}
        \begin{leftbar}[0.9\linewidth]
            Es wurde eine eigene ScreenController Klasse hinzugefügt.
        \end{leftbar}
    \item \textit{Methode editCurrentResolution(String, DropDownSelectionChangedEvent) hinzugefügt}
        \begin{leftbar}[0.9\linewidth]
            Wird aufgerufen sobald eine Auswahl im Auflösungs-DropDown getroffen wird.
        \end{leftbar}
    \item \textit{Methode editCurrentLanguage(String, DropDownSelectionChangedEvent) hinzugefügt}
        \begin{leftbar}[0.9\linewidth]
            Wird aufgerufen sobald eine Auswahl im Sprache-DropDown getroffen wird.
        \end{leftbar}
    \item \textit{Methode editCurrentKIDifficulty(String, DropDownSelectionChangedEvent) hinzugefügt}
        \begin{leftbar}[0.9\linewidth]
            Wird aufgerufen sobald eine Auswahl im KI-Schwiriegskeits-DropDown getroffen wird.
        \end{leftbar}
    \item \textit{Methode editCurrentQuality(String, DropDownSelectionChangedEvent) hinzugefügt}
        \begin{leftbar}[0.9\linewidth]
            Wird aufgerufen sobald eine Auswahl im Qualitäts-DropDown getroffen wird.
        \end{leftbar}
    \item \textit{Methode toggleFullscreen(String, CheckBoxStateChangedEvent) hinzugefügt}
        \begin{leftbar}[0.9\linewidth]
            Wird aufgerufen sobald eine Auswahl in der Vollbild-Checkbox getroffen wird.
        \end{leftbar}
    \item \textit{Methode showGraphicsTab() hinzugefügt}
        \begin{leftbar}[0.9\linewidth]
            Zeigt den Grafik Einstellungs-Tab an.
        \end{leftbar}
    \item \textit{Methode showInternTab() hinzugefügt}
        \begin{leftbar}[0.9\linewidth]
            Zeigt den Internen Einstellungs-Tab an.
        \end{leftbar}
    \item \textit{Methode resetSettings() hinzugefügt}
        \begin{leftbar}[0.9\linewidth]
            Setzt die Einstellungen zurück.
        \end{leftbar}
\end{itemize}

\subsubsection{LoadingScreen}
\begin{itemize}
    \item \textit{Implementiert nicht mehr ScreenController, erbt dafür von OwnScreenController}
        \begin{leftbar}[0.9\linewidth]
            Es wurde eine eigene ScreenController Klasse hinzugefügt.
        \end{leftbar}
    \item \textit{Folgende Methoden entfernt}
        \begin{itemize}
            \item openHUD() : HUD wird an anderer Stelle geöffnet.
            \item generateMap() : Funktionalität wude in startGame() verschoben.
            \item getMapBody() : Nicht mehr benötigt.
        \end{itemize}
\end{itemize}

\subsubsection{LoadingUtil}
\begin{itemize}
    \item \textit{Klasse hinzugefügt}
        \begin{leftbar}[0.9\linewidth]
            In dieser Util Klasse existieren Funktionen, die vom LoadingScreen aufgerufen 
            werden, um alle benötigten Sachen zu laden. Mittels Callbacks wird der 
            ladebalken entsprechend gesetzt.
        \end{leftbar}
\end{itemize}

\subsubsection{Hud}
\begin{itemize}
    \item \textit{Implementiert nicht mehr ScreenController, erbt dafür von OwnScreenController}
        \begin{leftbar}[0.9\linewidth]
            Es wurde eine eigene ScreenController Klasse hinzugefügt.
        \end{leftbar}
    \item \textit{Alle Funktionen entfernt/ersetzt}
        \begin{leftbar}[0.9\linewidth]
            Die Minimap wurde nicht im Hud integriert und die restlichen Funktionen wurden 
            ersetzt bzw entfernt. Funktionen um auf andere Bildschirme du wechseln wurde an 
            eine andere Stelle verlagert.
        \end{leftbar}
    \pagebreak
    \item \textit{Methode updateRank(int) hinzugefügt}
        \begin{leftbar}[0.9\linewidth]
            Aktualisiert den angezeigten Rang im Hud.
        \end{leftbar}
    \item \textit{Methode updateTime(float) hinzugefügt}
        \begin{leftbar}[0.9\linewidth]
            Aktualisiert die angezeigte Zeit im Hud.
        \end{leftbar}
    \item \textit{Methode updateCountdown(String) hinzugefügt}
        \begin{leftbar}[0.9\linewidth]
            Aktualisiert den angezeigten Countdown im Hud.
        \end{leftbar}
    \item \textit{Methode updateItem(String) hinzugefügt}
        \begin{leftbar}[0.9\linewidth]
            Aktualisiert das angezeigte Item im Hud.
        \end{leftbar}
\end{itemize}

\subsubsection{ErrorScreen}
\begin{itemize}
    \item \textit{Klasse hinzugefügt}
        \begin{leftbar}[0.9\linewidth]
            Diese Klasse dient dazu, eventuell auftretende Exceptions/Fehler während dem Spiel 
            anzeigen zu können mittels eines Popups.
        \end{leftbar}
\end{itemize}

\subsubsection{PausemenuScreen}
\begin{itemize}
    \item \textit{Implementiert nicht mehr ScreenController, erbt dafür von OwnScreenController}
        \begin{leftbar}[0.9\linewidth]
            Es wurde eine eigene ScreenController Klasse hinzugefügt.
        \end{leftbar}
\end{itemize}

\subsubsection{EndmenuScreen}
\begin{itemize}
    \item \textit{Implementiert nicht mehr ScreenController, erbt dafür von OwnScreenController}
        \begin{leftbar}[0.9\linewidth]
            Es wurde eine eigene ScreenController Klasse hinzugefügt.
        \end{leftbar}
    \item \textit{Methode updatePlayerEndTime() entfernt}
        \begin{leftbar}[0.9\linewidth]
            Ersetzt durch folgende Methode.
        \end{leftbar}
    \item \textit{Methode updatePlayerEndTimeAndColor(int, float, String) hinzugefügt}
        \begin{leftbar}[0.9\linewidth]
            Aktualisiert die Zeit und Farbe des Spielers.
        \end{leftbar}
\end{itemize}
	\pagebreak

	\section{Unit Test Überblick}
	\subsection{InGame - Grafik}

\subsubsection{ParticleBasicTest}

\begin{itemize}
    \item \textit{PartikelMesh wird erstellt}
        \begin{leftbar}[0.9\linewidth]
            VertexBuffer sollten mit korrekter Größe initialisiert und
            weiterverarbeitet werden werden.
        \end{leftbar}
    \item \textit{Grafische Tests}
        \begin{leftbar}[0.9\linewidth]
            Viele grafisch basierte tests
        \end{leftbar}
\end{itemize}
	\pagebreak
	\subsection{InGame - Interface}

\subsubsection{TestTick}
\begin{itemize}
    \item \textit{testAddRemoveProperty}
        \begin{leftbar}[0.9\linewidth]
            Tests adding and removing IProperties to/from a DynamicGameObject.
        \end{leftbar}

    \item \textit{testAddEvent}
        \begin{leftbar}[0.9\linewidth]
            Tests adding a DynamicGameObjectEvent to a DynamicGameObject.
        \end{leftbar}

    \item \textit{testSync}
        \begin{leftbar}[0.9\linewidth]
            Tests synchronizing one DynamicGameObject to another DynamicGameObject.
        \end{leftbar}
\end{itemize}

\subsubsection{TestDynamicGameObject}
\begin{itemize}
    \item \textit{testAddRemoveDGO}
        \begin{leftbar}[0.9\linewidth]
            Tests adding and removing DynamicGameObjects to/from a Tick.
        \end{leftbar}

    \item \textit{testAddEvent}
        \begin{leftbar}[0.9\linewidth]
            Tests adding a TickEvent to a Tick.
        \end{leftbar}

    \item \textit{testGetEventsAsStream}
        \begin{leftbar}[0.9\linewidth]
            Tests whether the Events in the Stream provided by the DynamicGameObject are correct.
        \end{leftbar}

    \item \textit{testSync}
        \begin{leftbar}[0.9\linewidth]
            Tests synchronizing one DynamicGameObject to another DynamicGameObject.
        \end{leftbar}
\end{itemize}

\subsubsection{TestTickProcessor}
TestTickProcessor is a Test-Application, loading a small scene and using the TickProcessor
to update changes.
	\pagebreak
	\subsection{Assets}

\subsubsection{TestJsonInfoLoader}
\begin{itemize}
    \item \textit{testLoadAssetPack}
        \begin{leftbar}[0.9\linewidth]
            Testet das Laden eines AssetPacks, welches durch eine .json Datei definiert ist.
        \end{leftbar}

    \item \textit{testLoadAssetInfo}
        \begin{leftbar}[0.9\linewidth]
            Testet das Laden einer AssetInfo, welche durch eine .json Datei definiert ist.
        \end{leftbar}
\end{itemize}

\subsubsection{TestResourceModuleMain}
Test-Anwendung startet eine JME3.2 Anwendung und läd mit hilfe des Asset-Moduls eine Szene.
	\pagebreak
	\subsection{Generierung}

\subsubsection{Testszenarien}

Diese Test-Anwendungen wurden im Laufe der Implemetierung für einzelne Teilmodule
erstellt und Teilweise wieder gelöscht, oder zusammengefasst (jedoch in älteren commits vorhanden).
In folgeneden Bereichen wurde ausgiebig getestet:

\begin{itemize}
    \item \textit{MapGeneration}
        \begin{leftbar}[0.9\linewidth]
            Zusammensetzung der Map.
        \end{leftbar}
\end{itemize}
\begin{itemize}
    \item \textit{MapBody}
        \begin{leftbar}[0.9\linewidth]
            Umwandlung der Road in das Roadmodel des anderen Teams und zurückrechnung der
            Positionen aus dem Roadmodel in globale Koordinaten.
        \end{leftbar}
\end{itemize}
\begin{itemize}
    \item \textit{DomeGeneration}
        \begin{leftbar}[0.9\linewidth]
            Generierung des Bodens, Triangulierung, Farbsetzung, Kuppel.
        \end{leftbar}
\end{itemize}
\begin{itemize}
    \item \textit{SphereRoad}
        \begin{leftbar}[0.9\linewidth]
            Generierung einer valide Strecke in einer Sphäre.
        \end{leftbar}
\end{itemize}
\begin{itemize}
    \item \textit{TunnelRoad}
        \begin{leftbar}[0.9\linewidth]
            Generierung einer Verbindungsstrecke für einen Tunnel.
        \end{leftbar}
\end{itemize}
\begin{itemize}
    \item \textit{RoomRoad}
        \begin{leftbar}[0.9\linewidth]
            Generierung einer Verbindungsstrecke für einen Raum.
        \end{leftbar}
\end{itemize}
	\pagebreak
	\subsection{GUI}

\subsubsection{TestMenuConfig}
\begin{itemize}
    \item \textit{testGetGraphicsConfig}
        \begin{leftbar}[0.9\linewidth]
            Testet den Getter bzw. das Auslesen der GraphicsConfig Datei.
        \end{leftbar}
    \item \textit{testSaveGraphicsConfig}
        \begin{leftbar}[0.9\linewidth]
            Testet den Setter bzw. das Speichern der GraphicsConfig Datei.
        \end{leftbar}
    \item \textit{testResetGraphicsConfig}
        \begin{leftbar}[0.9\linewidth]
            Testet die Zurücksetzung der GraphicsConfig Datei.
        \end{leftbar}
    \item \textit{testGetInternalGameConfig}
        \begin{leftbar}[0.9\linewidth]
            Testet den Getter bzw. das Auslesen der InternalGameConfig Datei.
        \end{leftbar}
    \item \textit{testSaveInternalGameConfig}
        \begin{leftbar}[0.9\linewidth]
            Testet den Setter bzw. das Speichern der InternalGameConfig Datei.
        \end{leftbar}
    \item \textit{testResetInternalGameConfig}
        \begin{leftbar}[0.9\linewidth]
            Testet die Zurücksetzung der InternalGameConfig Datei.
        \end{leftbar}
    \item \textit{testGetSeedConfig}
        \begin{leftbar}[0.9\linewidth]
            Testet den Getter bzw. das Auslesen der SeedConfig Datei.
        \end{leftbar}
    \item \textit{testSaveSeedConfig}
        \begin{leftbar}[0.9\linewidth]
            Testet den Setter bzw. das Speichern der SeedConfig Datei.
        \end{leftbar}
    \item \textit{testAddASeed}
        \begin{leftbar}[0.9\linewidth]
            Testet die Hinzufügung eines Seeds zur SeedConfig Datei.
        \end{leftbar}
    \item \textit{testResetSeedConfig}
        \begin{leftbar}[0.9\linewidth]
            Testet die Zurücksetzung der SeedConfig Datei.
        \end{leftbar}
\end{itemize}

\subsubsection{TestSplashScreen}
\begin{itemize}
    \item \textit{testConstructor}
        \begin{leftbar}[0.9\linewidth]
            Testet den Konstruktor.
        \end{leftbar}
\end{itemize}

\pagebreak
\subsubsection{TestNiftyAppState}
\begin{itemize}
    \item \textit{testConstructorWithOneOwnScreenController}
        \begin{leftbar}[0.9\linewidth]
            Testet den Konstruktor mit einem übergebenem OwnScreenController und die
            Initialisierung von Nifty.
        \end{leftbar}
    \item \textit{testConstructorWithTwoOwnScreenController}
        \begin{leftbar}[0.9\linewidth]
            Testet den Konstruktor mit zwei übergebenen OwnScreenController und die
            Initialisierung von Nifty.
        \end{leftbar}
    \item \textit{testCleanup}
        \begin{leftbar}[0.9\linewidth]
            Testet die Cleanup Methode.
        \end{leftbar}
\end{itemize}

\subsubsection{TestMainmenu}
\begin{itemize}
    \item \textit{testConstructor}
        \begin{leftbar}[0.9\linewidth]
            Testet den Konstruktor.
        \end{leftbar}
    \item \textit{testKeyEvent}
        \begin{leftbar}[0.9\linewidth]
            Testet ob ein KeyEvent korrekt behandelt wird.
        \end{leftbar}
\end{itemize}

\subsubsection{TestMenuInputMapping}
\begin{itemize}
    \item \textit{testConvert}
        \begin{leftbar}[0.9\linewidth]
            Testet für alle Eingaben, ob sie korrekt konvertiet werden.
        \end{leftbar}
\end{itemize}

\subsubsection{TestSinglePlayerScreen}
\begin{itemize}
    \item \textit{testConstructor}
        \begin{leftbar}[0.9\linewidth]
            Testet den Konstruktor.
        \end{leftbar}
    \item \textit{testEditCurrentSeed}
        \begin{leftbar}[0.9\linewidth]
            Testet das Seed DropDown Menü.
        \end{leftbar}
    \item \textit{testEditCurrentSeedNumber}
        \begin{leftbar}[0.9\linewidth]
            Testet die Eingabe der Seed Nummer im entsprechenden Feld.
        \end{leftbar}
    \item \textit{testEditKICount}
        \begin{leftbar}[0.9\linewidth]
            Testet die Eingabe der Anzahl von KI's.
        \end{leftbar}
    \pagebreak
    \item \textit{testSaveSeedName}
        \begin{leftbar}[0.9\linewidth]
            Testet die Änderung des Namen eines ausgewählten Seeds. 
        \end{leftbar}
    \item \textit{testEditSeednameVisibility}
        \begin{leftbar}[0.9\linewidth]
            Testet das Verhalten des Stift-Buttons rechts neben des Drop-Down Menüs.
        \end{leftbar}
    \item \textit{testSaveSeedNumber}
        \begin{leftbar}[0.9\linewidth]
            Testet die Speicherung der Seed Nummer.
        \end{leftbar}
    \item \textit{testOnStartScreen}
        \begin{leftbar}[0.9\linewidth]
            Testet die onStartScreen Methode, also ob der StartButton sichtbar ist.
        \end{leftbar}
\end{itemize}

\subsubsection{TestSettingsmenu}
\begin{itemize}
    \item \textit{testConstructor}
        \begin{leftbar}[0.9\linewidth]
            Testet den Konstruktor.
        \end{leftbar}
    \item \textit{testChangingResolution}
        \begin{leftbar}[0.9\linewidth]
            Testet die Änderung der Auflösung.
        \end{leftbar}
    \item \textit{testChangingLaguage}
        \begin{leftbar}[0.9\linewidth]
            Testet die Änderung der Sprache.
        \end{leftbar}
    \item \textit{testChangingKIDiff}
        \begin{leftbar}[0.9\linewidth]
            Testet die Änderung der KI Schwierigkeit.
        \end{leftbar}
    \item \textit{testChangingQuality}
        \begin{leftbar}[0.9\linewidth]
            Testet die Änderung der Qualität.
        \end{leftbar}
    \item \textit{testChangingFullscreen}
        \begin{leftbar}[0.9\linewidth]
            Testet die Änderung vom Vollbildmodus.
        \end{leftbar}
    \item \textit{showGraphicsTab}
        \begin{leftbar}[0.9\linewidth]
            Testet die Anzeige des Grafik Einstellungs-Tab.
        \end{leftbar}
    \item \textit{showInternTab}
        \begin{leftbar}[0.9\linewidth]
            Testet die Anzeige des internen Einstellungs-Tab.
        \end{leftbar}
    \item \textit{resetSettings}
        \begin{leftbar}[0.9\linewidth]
            Testet die Zurücksetzung der Einstellungen.
        \end{leftbar}
    \pagebreak
    \item \textit{testOnStartScreen}
        \begin{leftbar}[0.9\linewidth]
            Testet die onStartScreen Methode wenn der zuvorige Bildschirm das Pausemenü war.
        \end{leftbar}
\end{itemize}

\subsubsection{TestLoadingScreen}
\begin{itemize}
    \item \textit{testConstructor}
        \begin{leftbar}[0.9\linewidth]
            Testet den Konstruktor.
        \end{leftbar}
\end{itemize}

\subsubsection{TestHud}
\begin{itemize}
    \item \textit{testConstructor}
        \begin{leftbar}[0.9\linewidth]
            Testet den Konstruktor.
        \end{leftbar}
    \item \textit{testUpdateRank}
        \begin{leftbar}[0.9\linewidth]
            Testet die Anzeige des Rangs im Hud.
        \end{leftbar}
    \item \textit{testUpdateTime}
        \begin{leftbar}[0.9\linewidth]
            Testet die Anzeige der Zeit im Hud.
        \end{leftbar}
    \item \textit{testUpdateCountdown}
        \begin{leftbar}[0.9\linewidth]
            Testet die Anzeige des Countdowns im Hud.
        \end{leftbar}
    \item \textit{testUpdateItem}
        \begin{leftbar}[0.9\linewidth]
            Testet die Anzeige des Items im Hud.
        \end{leftbar}
\end{itemize}

\subsubsection{TestErrorScreen}
\begin{itemize}
    \item \textit{testConstructor}
        \begin{leftbar}[0.9\linewidth]
            Testet den Konstruktor.
        \end{leftbar}
\end{itemize}
\pagebreak
\subsubsection{TestPausemenu}
\begin{itemize}
    \item \textit{testConstructor}
        \begin{leftbar}[0.9\linewidth]
            Testet den Konstruktor.
        \end{leftbar}
    \item \textit{testSaveSeed}
        \begin{leftbar}[0.9\linewidth]
            Testet die Speicherung des aktuell gespielten Seeds.
        \end{leftbar}
    \item \textit{testKeyEvent}
        \begin{leftbar}[0.9\linewidth]
            Testet die korrekte Behandlung eines KeyEvents.
        \end{leftbar}
\end{itemize}

\subsubsection{TestEndmenuScreen}
\begin{itemize}
    \item \textit{testConstructor}
        \begin{leftbar}[0.9\linewidth]
            Testet den Konstruktor.
        \end{leftbar}
    \item \textit{testSaveSeed}
        \begin{leftbar}[0.9\linewidth]
            Testet die Speicherung des aktuell gespielten Seeds.
        \end{leftbar}
    \item \textit{testUpdatePlayerEndtimeandColor}
        \begin{leftbar}[0.9\linewidth]
            Testet die Anzeige der Zeit und Farbe eines Spielers.
        \end{leftbar}
\end{itemize}
	\pagebreak
	
	\section{Gantt}
	\subsection{Geplanter Ablauf}
	\begin{figure}[htbp]
		\centering
		\includegraphics[width=\linewidth]{./Bilder/PSE_Gantt.pdf}
		\caption{Gantt - Diagramm}
	\end{figure}

	\pagebreak

	\subsection{Tatsächlicher Ablauf}
	\begin{figure}[htbp]
		\centering
		\includegraphics[width=\linewidth]{./Bilder/GanttJonas.png}
		\caption{Jonas - Gantt - Diagramm}
	\end{figure}
	\begin{figure}[htbp]
		\centering
		\includegraphics[width=\linewidth]{./Bilder/GanttArtur.png}
		\caption{Artur - Gantt - Diagramm}
	\end{figure}
	\begin{figure}[htbp]
		\centering
		\includegraphics[width=\linewidth]{./Bilder/GanttFrederik.png}
		\caption{Frederik - Gantt - Diagramm}
	\end{figure}
	\begin{figure}[htbp]
		\centering
		\includegraphics[width=\linewidth]{./Bilder/Gantt_Generierung.png}
		\caption{Lukas und Sidney - Gantt - Diagramm}
	\end{figure}

\end{document}